\documentclass[a4paper,11pt]{article}
\usepackage{iftex}
\ifPDFTeX
  \usepackage[T2A]{fontenc}
  \usepackage[utf8]{inputenc}
  \usepackage[russian]{babel}
\else
  \usepackage{fontspec}
  \setmainfont{Noto Serif}
  \usepackage[russian]{babel}
\fi
\usepackage[a4paper,margin=24mm,headheight=15pt]{geometry}
\usepackage{amsmath,amssymb,booktabs,tabularx,xcolor,tikz,fancyhdr,hyperref}
\definecolor{accent}{HTML}{245D60}
\hypersetup{colorlinks=true,linkcolor=accent,pdftitle={Ревью домашней работы — квадратные корни}}
\setlength{\parindent}{0pt}
\setlength{\parskip}{5pt}
\setlength{\abovedisplayskip}{5pt}
\setlength{\belowdisplayskip}{5pt}
\pagestyle{fancy}
\fancyhf{}
\fancyhead[L]{\small Квадратные корни}
\fancyhead[R]{\small Екатерина Скелли}
\fancyfoot[C]{\thepage}
\renewcommand{\headrulewidth}{0.3pt}
\newcommand{\err}{\textcolor{red!75!black}{\textbf{Ошибка:}}}
\newcommand{\idea}{\textcolor{blue!60!black}{\textbf{Ключевая идея:}}}
\newcommand{\result}{\textcolor{teal!60!black}{\textbf{Итог:}}}
\newcommand{\taskhead}[1]{\phantomsection\label{task:#1}\section*{\textcolor{accent}{Задание № #1}}}
\begin{document}
\begin{titlepage}
\centering
\vspace*{22mm}
{\large МАТЕМАТИКА\par}
\vspace{24mm}
{\LARGE\bfseries Ревью домашней работы\par}
\vspace{7mm}
{\Large Квадратные корни\par}
\vspace{18mm}
{\large Екатерина Скелли\par}
\vspace{5mm}
28 ноября 2026 г.\par
\vfill
Лёвин Артём Александрович\par
эксклюзивно для Екатерины Скелли\par
\vspace{10mm}
\begin{tikzpicture}
\draw[accent,line width=.7pt] (0,0) .. controls (2.1,0.65) and (4.4,-0.65) .. (7,0);
\end{tikzpicture}
\end{titlepage}
\setcounter{page}{1}
\section*{\textcolor{accent}{Сводная таблица}}
\small
\renewcommand{\arraystretch}{1.45}
\begin{tabularx}{\linewidth}{@{}c p{25mm} X p{23mm} p{39mm}@{}}
\toprule
№ & Тема & Суть ошибки & Ваш ответ & Правильный ответ\\
\midrule
\hyperref[task:8]{8} & Разность дробей & При приведении к общему знаменателю вычитание заменено сложением. & $4$ & $-2\sqrt3$\\
\hyperref[task:10]{10} & Корень, модуль, ОДЗ & Перепутаны условия раскрытия модуля; решение не завершено. & Не указан & $\begin{cases}3-2x,&-1\le x<\frac23,\\4x-1,&x\ge\frac23.\end{cases}$\\
\bottomrule
\end{tabularx}
\normalsize
\clearpage
\thispagestyle{empty}
\vspace*{\fill}
\begin{center}
{\LARGE\bfseries\textcolor{accent}{Разбор ошибок}}
\end{center}
\vspace*{\fill}
\clearpage
\taskhead{8}
\textbf{Текст задания.} Вычислите
\[
E=\frac{\sqrt3-1}{\sqrt3+1}-\frac{\sqrt3+1}{\sqrt3-1}.
\]
\textbf{Ваше решение.} Вы выбрали общий знаменатель $(\sqrt3+1)(\sqrt3-1)$, затем записали сумму квадратов в числителе:
\[
\frac{(\sqrt3-1)^2+(\sqrt3+1)^2}{(\sqrt3+1)(\sqrt3-1)}
=\frac{8}{2}=4.
\]
\textbf{Ваш ответ:} $4$.

\err\ Потеря знака вычитания при приведении дробей к общему знаменателю.

\textbf{Где именно возникает ошибка.} Выбор знаменателя верен. Первый неверный переход --- запись знака плюс между квадратами в числителе. Исходная операция вычитания должна сохраниться.

\textbf{Почему это неверно.} При ненулевых $b$ и $d$ действует правило
\[
\frac ab-\frac cd=\frac{ad-bc}{bd}.
\]
Вычитается всё произведение $bc$. После раскрытия скобок знак минус перед вторым квадратом меняет знак каждого его слагаемого.

\idea\ Сопряжённые выражения дают разность квадратов в знаменателе. В числителе также возникает разность квадратов; знак исходной операции определяет результат.

\textbf{Исправление и полное решение.} Оба знаменателя ненулевые, поскольку $\sqrt3>1$. Сохраняем выбранный Вами общий знаменатель:
\begin{align*}
E&=\frac{(\sqrt3-1)^2-(\sqrt3+1)^2}{(\sqrt3+1)(\sqrt3-1)}\\
&=\frac{(3-2\sqrt3+1)-(3+2\sqrt3+1)}{3-1}\\
&=\frac{3-2\sqrt3+1-3-2\sqrt3-1}{2}\\
&=\frac{-4\sqrt3}{2}=-2\sqrt3.
\end{align*}
\textbf{Проверка.} Первая исходная дробь находится между нулём и единицей, вторая больше единицы. Их разность отрицательна; положительный ответ $4$ этому противоречит.

\result\ $\boxed{-2\sqrt3}$.

\textbf{Задача на закрепление.} Вычислите
\[
\frac{\sqrt5-2}{\sqrt5+2}-\frac{\sqrt5+2}{\sqrt5-2}.
\]
\clearpage
\taskhead{10}
\textbf{Текст задания.} Упростите выражение, укажите ОДЗ и запишите ответ без модуля по случаям:
\[
F(x)=\sqrt{(3x-2)^2}+\bigl(\sqrt{x+1}\bigr)^2.
\]
\textbf{Ваше решение.} Верно использовано равенство $\sqrt{(3x-2)^2}=|3x-2|$. В записи случаев условие $x\le\frac23$ поставлено рядом с $3x-2$, а условие $x\ge\frac23$ --- рядом с $-3x+2$. В рукописи видны пометки около этих условий; окончательное исправленное выражение отсутствует.

\textbf{Ваш ответ:} не указан.

\err\ Условия раскрытия модуля перепутаны в записанных строках. Кроме того, не указана ОДЗ и не учтено второе слагаемое в окончательном результате.

\textbf{Где именно возникает ошибка.} Последний правильный шаг: $\sqrt{(3x-2)^2}=|3x-2|$. Далее условия должны определяться знаком выражения $3x-2$.

\textbf{Почему это неверно.} Арифметический квадратный корень всегда неотрицателен. Для любого действительного $u$ выполняется
\[
\sqrt{u^2}=|u|=\begin{cases}u,&u\ge0,\\-u,&u<0.\end{cases}
\]
Например, при $x=0$ выражение $3x-2=-2$, поэтому его модуль равен $2$. Если оставить $3x-2$ на промежутке $x<2/3$, получится отрицательное значение модуля.

\idea\ Сначала найдите ОДЗ всего выражения. Затем определите знак подмодульного выражения, раскройте модуль и добавьте второе слагаемое в каждом случае.

\textbf{Исправление и полное решение.}

\textbf{Шаг 1. ОДЗ.} Квадрат $(3x-2)^2\ge0$ при любом действительном $x$. Второй корень существует при $x+1\ge0$, откуда
\[
\boxed{x\in[-1;+\infty)}.
\]
\textbf{Шаг 2. Упрощение корней.} На найденной ОДЗ $(\sqrt{x+1})^2=x+1$, следовательно,
\[
F(x)=|3x-2|+x+1.
\]
\textbf{Шаг 3. Определение границы случаев.}
\[
3x-2\ge0\iff x\ge\frac23,\qquad
3x-2<0\iff x<\frac23.
\]
Поэтому $|3x-2|=3x-2$ справа от границы (включая её), а слева $|3x-2|=-(3x-2)=2-3x$.
\clearpage
\section*{\textcolor{accent}{Задание № 10. Продолжение}}
\textbf{Шаг 4. Пересечение с ОДЗ и сложение.}

Первый случай: $-1\le x<\frac23$. Подмодульное выражение отрицательно, поэтому
\[
F(x)=2-3x+x+1=3-2x.
\]
Второй случай: $x\ge\frac23$. Подмодульное выражение неотрицательно, поэтому
\[
F(x)=3x-2+x+1=4x-1.
\]
\begin{center}
\begin{tikzpicture}[x=2.2cm,y=1cm]
\draw[->,thin] (-1.55,0)--(2.2,0) node[right] {$x$};
\draw[accent,line width=1.5pt] (-1,0)--(.66667,0);
\draw[blue!55!black,line width=1.5pt,->] (.66667,0)--(2.05,0);
\fill[accent] (-1,0) circle (2pt) node[below=4pt,black] {$-1$};
\fill[blue!55!black] (.66667,0) circle (2pt) node[below=4pt,black] {$\frac23$};
\node[above=5pt,accent] at (-.17,0) {$3-2x$};
\node[above=5pt,blue!55!black] at (1.45,0) {$4x-1$};
\end{tikzpicture}
\end{center}
\textbf{Проверка.} Возьмём по одному значению из каждого промежутка и отдельно проверим границу.
\begin{align*}
F(0)&=\sqrt{(-2)^2}+(\sqrt1)^2=2+1=3,\\
F(1)&=\sqrt{1^2}+(\sqrt2)^2=1+2=3,\\
F\left(\frac23\right)&=0+\frac23+1=\frac53.
\end{align*}
По полученным формулам: $3-2\cdot0=3$, $4\cdot1-1=3$, $4\cdot\frac23-1=\frac53$. На левой границе ОДЗ $F(-1)=5$; обе части исходного выражения определены.

\result\ ОДЗ: $x\in[-1;+\infty)$;
\[
\boxed{F(x)=\begin{cases}3-2x,&-1\le x<\frac23,\\[2pt]4x-1,&x\ge\frac23.\end{cases}}
\]
\textbf{Задача на закрепление.} Упростите выражение, укажите ОДЗ и запишите ответ без модуля по случаям:
\[
\sqrt{(2x-5)^2}+\bigl(\sqrt{x+3}\bigr)^2.
\]
\vfill
\begin{center}\small Лёвин Артём Александрович\par эксклюзивно для Екатерины Скелли\end{center}
\end{document}
